\section{Experiments}

\noindent \textbf{Datasets.} We used the following brain MRI datasets for training and evaluation (\# of \textcolor{RoyalBlue}{train} / \textcolor{ForestGreen}{validation} / \textcolor{Orange}{test} subjects): BraTS2021 \cite{baid2021rsna} (\textcolor{RoyalBlue}{1101} / \textcolor{ForestGreen}{40} / \textcolor{Orange}{110}), BraTS2023-MEN \cite{labella2023asnrmiccai} (\textcolor{RoyalBlue}{880} / \textcolor{ForestGreen}{29} / \textcolor{Orange}{91}), IXI \footnote{https://brain-development.org/ixi-dataset/} (\textcolor{RoyalBlue}{500} / \textcolor{ForestGreen}{27} / \textcolor{Orange}{51}), and ATLAS 2.0 \cite{liew2022large} (\textcolor{RoyalBlue}{430} / \textcolor{ForestGreen}{23} / \textcolor{Orange}{43}). BraTS2021 and BraTS2023-MEN include T1, T1ce, T2, and FLAIR modalities; IXI contains T1, T2, and PD modalities; and ATLAS 2.0 consists of the T1 modality. We also employed UPENN-GBM \cite{bakas2021multi} (\textcolor{Orange}{147}) for an external test dataset (image-to-image tasks only).

\vspace{3pt}
\noindent \textbf{Implementation Details.} We used \texttt{dolly-v2-3b} \cite{conover2023free} as our backbone language model, which has a vocabulary size of 50,821 tokens ($K_\texttt{text}$). For the image tokenizer, we used VQGAN \cite{esser2021taming} with a codebook $K_\texttt{image}$ of size 1024. By adding 1024 image tokens, we extended the embedding table to a total of 51,845 tokens ($K_\texttt{text} + K_\texttt{image}$). We used ROUGE \cite{lin-2004-rouge}, METEOR \cite{banerjee-lavie-2005-meteor}, and BERT-F1 \cite{Zhang2020BERTScore} to evaluate report generation. For segmentation, we measured performance using the Dice score. For image translation, we used PSNR, SSIM, and FID \cite{heusel2017gans} as metrics. The model was trained on four A100 80GB GPUs with a batch size of 2 per GPU, a learning rate of 5e-6, and the AdamW \cite{loshchilov2018decoupled} optimizer. Additional implementation details are provided in the Appendix.

